Thus $C,C'$ are two Cartan subgroups of $H$ (by a)) which contain the
same regular element $g$ of $H$, and are therefore equal. \hfill Q.E.D.

Let us prove c): denoting by $\nu(u)$ the nullity\nde{That is, the
dimension of its nilspace.} of $\mathrm{id}-u$, for an endomorphism of
a finite-dimensional vector space, one will have
\[
 \nu(\ad(g)_{\mathfrak g})=\nu(\ad(g)_{\mathfrak h})
 +\nu(\ad(g)_{\mathfrak g/\mathfrak h}),
\]
now the two terms on the right-hand side are respectively $\geq$ the
nilpotent rank of $H$ (equal to the nilpotent rank $\rho$ of $G$, by a))
and $\geq0$, so one has $\nu(\ad(g)_{\mathfrak g})=\rho$ if and only
if $\nu(\ad(g)_{\mathfrak h})=\rho$ and
$\nu(\ad(g)_{\mathfrak g/\mathfrak h})=0$, i.e. $g$ is regular in $G$
if and only if $g$ is regular in $H$ and $\ad(g)_{\mathfrak g/\mathfrak h}$
has no nonzero invariants.
\end{proof}

\begin{remarks}{2.9}\label{XIII.2.9}
In statement 2.1, one cannot weaken condition (iii) by supposing merely
that $H$ contains a maximal torus and has finite index in its normalizer,
even if one requires that this normalizer moreover be smooth, i.e. that
one have $H=N^0$, and even when $G$ is affine and solvable. An example
is provided by the group $G$ of matrices of the form
\[
 g=\begin{pmatrix}t&a&c\\0&1&b\\0&0&1\end{pmatrix},
\]
and the subgroup $H$ of matrices of the preceding form, with $b=c=0$
(N.B. the Cartan subgroup of $G$ here consists of the matrices $g$
with $a=c=0$).
\end{remarks}

\begin{remarks}{2.10}\label{XIII.2.10}
Let $G$ be a smooth algebraic group over $k$, $H$ a smooth algebraic
subgroup, but no longer suppose that $H$ and $G$ are connected. Suppose
that $H^0$ contains a Cartan subgroup of $G^0$. Then
$(\mathfrak g/\mathfrak h)^H\subset(\mathfrak g/\mathfrak h)^{H^0}=0$,
so $H^0=N^0$ ($N$ is the normalizer of $H$); in particular $N$ is smooth.
However, one easily constructs examples, with $G$ connected, where $H$
has a connected component $H_i$ such that for no $h\in H_i(k)$ does
one have $(\mathfrak g/\mathfrak h)^{\ad(h)}=0$, i.e. the morphism
$(g,h)\mto\ad(g)\cdot h$ from $G\times H_i$ to $G$ is \'etale
(or even quasi-finite) at no point (take, for example, for $H$ the
normalizer of the maximal torus in $\SL(2)_k$). Likewise, even if the
image $H'$ of $H$ in the finite group $G'=G/G^0$ is equal to $G'$, it
is not necessarily true that the union of the conjugates of $H$ in
$G$ is dense (take for $G$ the semidirect product of $\ZZ/2\ZZ$ with
$G^0=\SL(2)_k$, on which it acts by ``symmetry'', and for $H$ the
semidirect product $\ZZ/2\ZZ\cdot T$, where $T$ is a maximal torus
of $G^0$). On the other hand, if one does not assume a priori that
$H^0$ contains a Cartan subgroup of $G^0$, but that the union of the
conjugates of $H$ in $G$ is dense, then $H^0$ necessarily contains a
Cartan subgroup of $G^0$: to verify this, one may clearly suppose
$G$ connected, and it suffices to repeat the proof of 2.1 (vi)
$\Rightarrow$ (i), which is valid without supposing $H$ connected.
% typo: closing parenthesis added after the SL(2)_k example.
\end{remarks}

\sgasection{3}{Case of an arbitrary base prescheme}
Suppose first that one is over a base field $k$, not necessarily
algebraically closed. Since conditions 2.1 (i bis), (iv bis), (v),
(vi) are invariant under extension of the base field, one sees by
passing to the algebraic closure $\bar k$ of $k$ that they are equivalent
to one another, and equivalent to the fact that $H_{\bar k}$ contains
a Cartan subgroup of $G_{\bar k}$. When this condition is satisfied,
then (with the notation of 2.3) it will still be true that $k(X)$ is
a finite separable extension of $k(G)$, of degree $d$ independent of
any extension of the base. If $U$ is the largest open subset of $G$
such that $\psi$ induces a morphism $\psi^{-1}(U)\to U$ which is finite
and \'etale, then the formation of $U$ commutes with extension of the
base field. If $\psi$ is birational, then $U$ is also the largest open
subset of $G$ such that $\psi$ induces an isomorphism
$\psi^{-1}(U)\isomto U$, and if then $g\in U(k)$, there exists one and
only one subgroup $H'$ of $G$, conjugate to $H$ over the algebraic
closure $\bar k$ of $k$, such that $H'\ni g$.

A point $g\in G(k)$ is said to be \emph{regular} if it is regular as an
element of $G(\bar k)=G_{\bar k}(\bar k)$. More generally, the construction
of 2.7 gives one an open subset of $G$, whose formation commutes with
every extension of the base field, called the \emph{open subset of
regular points} of $G$, which is also characterized by the fact that
for every algebraically closed extension $K$ of $k$ and every point
$g\in G(K)$, $g$ is a regular point of $G_K$ if and only if $g\in U(K)$.
If $g\in U(k)$\nde{Correction of $G(k)$ to $U(k)$.}, one sees that $g$
is contained in one and only one Cartan subgroup of $G$, as we shall
show under more general conditions below.
% typo: on ouvert -> un ouvert; voit -> on voit.

Let $G$ be a smooth separated group prescheme of finite type with
connected fibers over the prescheme $S$; consider the functor
$\mathcal C:(\Sch)^\circ_{/S}\to\Ens$ defined by
\[
 \mathcal C(S')=\text{set of Cartan subgroups (XII 3.1) of }G_{S'}.
\]
Suppose that this functor is representable by a prescheme smooth over
$S$; we give in XV a necessary and sufficient condition for this to
be so, but we already know that this hypothesis is satisfied if $G$
is affine over $S$ of locally constant reductive rank (XII 3.3), or
more generally if $G$ admits a maximal torus locally for the fpqc
topology (XII 7.1 a)), for example if $S$ is the spectrum of a field.
Let $X$ be the Cartan subgroup of the $\mathcal C$-group prescheme
$G_{\mathcal C}$, the ``universal Cartan subgroup'' of $G$. As a
prescheme over $S$, $X$ therefore represents the functor
\[
\begin{aligned}
 X(S')=\{&(C,g): C\text{ a Cartan subgroup of }G_{S'}\text{ and }g\text{ a}\\
 &\text{section of }C\text{ over }S'\}.
\end{aligned}
\]
Consider the canonical projection morphism $(C,g)\mto g$
\[
 \psi:X\to G.
\]
One then has the

\begin{theorem}{3.1}\label{XIII.3.1}
Under the preceding conditions on $G$, and with the preceding notation,
let $U$ be the set of $g\in G$ such that $g$ is a regular element of
its fiber $G_s$. Then $U$ is open, and it is also the largest open subset
$U$ of $G$ such that $\psi$ induces an isomorphism
\[
 \psi^{-1}(U)\isomto U.
\]
\end{theorem}
\begin{proof}
Let us first prove that $U$ is open. From the hypothesis of representability
of $\mathcal C$ as a prescheme smooth over $S$, since its structural
morphism is clearly surjective, one concludes at once that $G$ admits a
Cartan subgroup locally for the \'etale topology, and that the nilpotent
rank of the fibers of $G$ is locally constant. The same holds for the
dimension of the fibers of $G$, and, localizing on $S$, one may suppose
that both are constant, say $\rho$ and $n$. Consider then the Killing
polynomial
\[
 P_G(t)=t^n+c_1t^{n-1}+\cdots+c_n\in A[t],
 \qquad\text{where }A=\Gamma(G,\OO_G).
\]
The restriction of this polynomial to the fibers $G_s$ of $G$, and in
particular to the fibers at the maximal points of $S$, is divisible
by $(t-1)^\rho$, which is expressed by the fact that certain linear
combinations with integer coefficients of the $c_i$ are zero on the
fibers $G_s$. When $S$ is reduced (which one may suppose in order to
establish that $U$ is open), it follows that they are themselves zero,
hence that the Killing polynomial itself is divisible by $(t-1)^\rho$,
say
\[
 P_G(t)=(t-1)^\rho(t^{n-\rho}+b_1t^{n-r-1}+\cdots+b_{n-\rho}).
\]
% typo?: kept as printed: n-r-1, with r rather than rho.
Let $b$ be the sum of the coefficients $b_0=1,b_1,\ldots,b_{n-\rho}$ of
the second factor; then by 2.7 applied to the fibers of $G$, one sees that
\[
 U=G_b,
\]
which indeed proves that $U$ is open.

To prove that $\psi^{-1}(U)\to U$ is an isomorphism, one is reduced by
SGA1 I 5.7 to verifying it fiber by fiber, which reduces one to the case
of a base field, which may be supposed algebraically closed. Then there
exists a Cartan subgroup $C$ of $G$, and if $N$ is its normalizer,
$\mathcal C$ is identified by the conjugacy theorem XII 7.1 a) and b)
with $G/N$, and the morphism $\psi:X\to G$ considered here is precisely
that defined in \No 2. One then concludes by 2.6 b). The same argument
also shows that $U$ is the largest open subset of $G$ such that $\psi$
induces an isomorphism $\psi^{-1}(U)\to U$.
\end{proof}

\begin{corollary}{3.2}\label{XIII.3.2}
Under the conditions of 3.1, let $g$ be a regular section of $G$, i.e.
such that for every $s\in S$, $g(s)$ is a regular point of $G_s$. Then
there exists one and only one Cartan subgroup $C$ of $G$ such that $g$
is a section of $C$.
\end{corollary}
Indeed, the hypothesis on $g$ means that $g$ is a section of $U$, and
the conclusion that there exists a unique section of $X$ lifting it,
which is merely another way of expressing that $\psi^{-1}(U)\to U$
is an isomorphism.

Note now that the open subset $\psi^{-1}(U)$ of the Cartan subgroup
$X$ of $G_{\mathcal C}$ is precisely the open subset of $X$ consisting
of the points of $X$ which are regular in $G_{\mathcal C}$ (by which
one means: regular in their fiber). One thus obtains a natural
``fibration'' of the dense open subset $U$ of regular points of $G$
over the prescheme $\mathcal C$, the fibers being dense open subsets
of Cartan subgroups of the fibers of $G_{\mathcal C}$ (namely the
open subsets of regular points in $G$). One finds, for example, the
following result (which will be made considerably more precise in
the next expos\'e):

\begin{corollary}{3.3}\label{XIII.3.3}
Let $G$ be a smooth connected algebraic group over the field $k$,
$\mathcal T$ the scheme of maximal tori of $G$ ($\simeq$ the scheme
of Cartan subgroups of $G$). Then the function field $k(G)$ of $G$
is isomorphic to the function field of a smooth connected affine
nilpotent algebraic group $C$ over the function field $k(\mathcal T)$
of $\mathcal T$, namely $C=$ ``the generic Cartan subgroup of $G$''.
If $G$ is affine of unipotent rank zero, i.e. if the Cartan subgroups
of $G_K$ are tori, then $k(G)$ is a unirational extension of
$k(\mathcal T)$.
\end{corollary}
Of course, by the generic Cartan subgroup of $G$, one means (by abuse
of language) the Cartan subgroup of $G_{k(\mathcal T)}$ which is the
generic fiber of $X$ over $\mathcal T$. It only remains to prove the
last assertion of 3.3, which is contained in the following well-known
result (due to Chevalley):

\begin{lemma}{3.4}\label{XIII.3.4}
Let $k$ be a field, $T$ a torus over $k$, and $k(T)$ the field of rational
functions on $T$; then $k(T)$ is a unirational extension of $k$, i.e.
is contained in a purely transcendental extension of $k$.
\end{lemma}
Indeed, let $k'$ be a finite separable extension of $k$ which splits\nde{Splits
(?), or: ``over which $T$ is split''.} $T$ (X 1.4); then $T\otimes_k k'$
is a rational variety, i.e. admits a dense open subset isomorphic to a
dense open subset of affine space $\mathbb A^n_{k'}$, hence
$T'=\prod_{\Spec(k')/\Spec(k)}T_{k'}/\Spec(k')$ is a rational variety
(for it admits a dense open subset isomorphic to a dense open subset
of $\prod_{\Spec(k')/\Spec(k)}\mathbb A^n_{k'}$, which is isomorphic to
affine space of dimension $mn$ over $k$, where $m=[k':k]$). Consider
the norm homomorphism from $T'$ to $T$ (defined as soon as $T$ is a
commutative group scheme over $k$); the composite $T\to T'\to T$ is
the $m$th power on $T$, hence dominant, so $T'\to T$ is dominant,
which proves that $T$ is unirational.

Let us return to the conditions of 3.1, but supposing even that $G$
admits a maximal torus locally for the fpqc topology (XII 7.1). Let
$Y$ be the maximal torus of the Cartan subgroup $X$ of $G$, so that
the morphism $\psi:X\to G$ induces a morphism $Y\to G$ whose image
consists, as a set, of the semisimple elements of the fibers of $G$
(XII 8). Finally, it follows from 3.1 that the restriction of $\psi$
to the open subset $Y^{\mathrm{reg}}$ of regular points of $Y$ induces
a closed immersion
\[
 Y^{\mathrm{reg}}\to U=G^{\mathrm{reg}}.
\]
Making explicit the meaning of $Z=Y^{\mathrm{reg}}$ as a functor over
$S$, one finds:
\begin{corollary}{3.5}\label{XIII.3.5}
Let $G$ be a smooth separated $S$-group prescheme of finite type with
connected fibers over the prescheme $S$, admitting a maximal torus
locally for the fpqc topology. Let $Z:(\Sch)^\circ_{/S}\to\Ens$ be the
functor defined by
\[
\begin{aligned}
 Z(S')=\{&\text{regular sections of }G_{S'}\text{ over }S'\text{ which are
 contained in}\\
 &\text{a maximal torus of }G_{S'}\}.
\end{aligned}
\]
Then $Z$ is representable by a closed subprescheme, smooth over $S$,
of the open subset $U=G^{\mathrm{reg}}$ of $G$ introduced in 3.1.
\end{corollary}
Finally, note the following result, which makes the density theorem
2.1 (i) $\Rightarrow$ (vi) more precise:

\begin{corollary}{3.6}\label{XIII.3.6}
Under the conditions of 3.5, let $C$ be a Cartan subgroup of $G$,
and consider the morphism
\[
 \varphi:Z\times C\to G
\]
defined by $\varphi(g,h)=\ad(g)h=ghg^{-1}$. Then $\varphi$ is dominant.
\end{corollary}
It clearly suffices to prove this fiber by fiber, which reduces one
to the case where $S$ is the spectrum of an algebraically closed field.
Let $T$ be the maximal torus of $C$, $t_0$ an element of $T(k)$ regular
in $G$, and $c_0$ an element of $C(k)$ regular in $G$; consider
$\varphi^{-1}(\varphi(t_0,c_0))$, whose points rational over $k$ are the
pairs $(t,c)$, with $t\in Z(k)$, $c\in C(k)$, such that
\[
 \ad(t)c=\ad(t_0)c_0\qquad\text{i.e.}\qquad c=\ad(t^{-1}t_0)c_0,
\]
which are therefore in one-to-one correspondence with the $t\in Z(k)$
such that $\ad(t^{-1}t_0)c_0\in C$, or, what amounts to the same thing,
$c_0$ being regular, such that $t^{-1}t_0\in N$ (the normalizer of $C$),
i.e. $t\in N$. One obtains an open and closed part of this fiber by
restricting to the $t\in Z(k)$ such that $t\in C(k)$. Thus one has
found a connected component of $\varphi^{-1}\varphi(t_0,c_0)$ isomorphic
to $T$ (N.B. that the preceding set-theoretic argument does indeed give
an isomorphism of schemes is seen by replacing points with values in
$k$ by points with values in an arbitrary $k$-prescheme), so the generic
fiber of $\varphi$ has dimension $\leq\dim T$, hence $\Im(\varphi)$ has
dimension $\geq\dim Z\times C-\dim T=\dim Z+\dim C-\dim T$; now one has
$\dim Z=\dim Y=\dim\mathcal C+\dim T=\dim G-\dim C+\dim T$, whence
finally $\dim\Im(\varphi)\geq\dim G$, so $\varphi$ is dominant.
\hfill Q.E.D.

\begin{remarks}{3.7}\label{XIII.3.7}
One will note that the argument shows in addition that the connected
component at $(t_0,c_0)$ of the fiber of $\varphi(t_0,c_0)$ is isomorphic
to $T$, in particular is smooth over $k$, and has the same dimension
as the generic fiber, which implies that $\varphi$ is in fact smooth
at $(t_0,c_0)$ (which one should also be able to verify by calculating
the tangent map). It follows that under the conditions of 3.6 the
induced morphism $Z\times_S C^{\mathrm{reg}}\to G^{\mathrm{reg}}$ (where
one has set $C^{\mathrm{reg}}=C\cap G^{\mathrm{reg}}$) is a smooth morphism.
One sees likewise that the analogous morphism
$Z\times T^{\mathrm{reg}}\to Z$ (where $T$ is a maximal torus of $G$)
is smooth; more generally, for every smooth connected invariant
algebraic subgroup $H$ of $C$ containing a regular element $c_0$ of
$G(k)$, the image of $Z\times H\to G$ is dense in that of $G\times H\to G$.
\end{remarks}

\sgasection{4}{Lie algebras over a field: rank, regular elements, Cartan subalgebras}
In the remainder of this expos\'e, we take up again the theory developed
by Chevalley in his book ``Th\'eorie des Groupes de Lie III'' (Act. Sc.
Ind. 1226, Paris 1955), the technique of schemes allowing us to eliminate
the characteristic-zero hypothesis. We begin by recalling in this
section certain well-known notions and results.

Let $\mathfrak g$ be a Lie algebra over a ring $k$. For every
$a\in\mathfrak g$, one denotes by $\ad(a)$ the endomorphism
\[
 \ad(a)\cdot x=[a,x]
\]
of $\mathfrak g$, which is a derivation of the Lie algebra $\mathfrak g$.
Now, for every derivation $D$ of the Lie algebra $\mathfrak g$, the
nilspace of $D$, i.e. the union of the kernels of the iterates of $D$,
is a Lie subalgebra of $\mathfrak g$, as one sees from Leibniz's formula
\[
 D^m([x,y])=\sum_{0\leq p\leq m}\binom{m}{p}[D^px,D^{m-p}y].
\]
We shall set
\[
 \operatorname{Nil}(a,\mathfrak g)=\text{nilspace of }\ad(a)
 =\bigcup_{m\geq0}\Ker\ad(a)^m;
\]
when no confusion is to be feared, we shall denote it simply by
$\operatorname{Nil}(a)$, and call it the nilspace of $a$ (in $\mathfrak g$).

\begin{proposition}{4.1}\label{XIII.4.1}
For every $a\in\mathfrak g$, its nilspace $\operatorname{Nil}(a,\mathfrak g)$
is a Lie subalgebra of $\mathfrak g$, equal to its own normalizer.
\end{proposition}
It remains to prove that it is its own normalizer, i.e. that every
element of $\mathfrak g/\operatorname{Nil}$ annihilated by the adjoint
representation of $\operatorname{Nil}$ in
$\mathfrak g/\operatorname{Nil}$ is zero, which is trivial (for every
element in this quotient annihilated by $\Ad(a)$ is zero).
% typo?: kept as printed: capital Ad(a).

In the rest of this section, we suppose that $k$ is a field, and
$\mathfrak g$ finite-dimensional over $k$. We shall denote by
$\mathrm W(\mathfrak g)$ the scheme over $k$ defined by $\mathfrak g$,
whose points in the $k$-algebra $A$ are the elements of
$\mathfrak g\otimes_k A$. If $a\in\mathfrak g$, the characteristic
polynomial of $\ad(a)$ is also called the characteristic polynomial
or \emph{Killing polynomial} of $a$ in $\mathfrak g$, say
\[
 P_{\mathfrak g}(a,t)=t^n+c_1(a)t^{n-1}+\cdots+c_n(a),
\]
where $n=\operatorname{rank}_k\mathfrak g$, the $c_i(a)\in k$. Taking
this polynomial also for $a\in\mathfrak g\otimes_k A$\nde{$\mathrm W\otimes_k A$
has been corrected to $\mathfrak g\otimes_k A$.}, where $A$ is an
arbitrary $k$-algebra, one sees that the $c_i(a)$ come from well-determined
sections $c_i$ of the structure sheaf of $\mathrm W(\mathfrak g)$, i.e.
elements of the symmetric algebra $A=\operatorname{Sym}_k(\mathfrak g^\vee)$,
where $\mathfrak g^\vee$ is the dual of the $k$-module $\mathfrak g$.
(When $k$ is an infinite field, the $c_i$ are determined by knowledge
of the corresponding polynomial functions $\mathfrak g\to k$, but
this is no longer so if $k$ is a finite field.) Let $r$ be the largest
integer such that the Killing polynomial
\[
 P_{\mathfrak g}(t)=t^n+c_1t^{n-1}+\cdots+c_n\in A[t]
\]
is divisible by $t^r$, i.e. one has:
\[
 P_{\mathfrak g}(t)=t^n+c_1t^{n-1}+\ldots+c_{n-r}t^r,
 \qquad c_{n-r}\ne0.
\]
The integer $r$ is called the \emph{nilpotent rank} of the Lie algebra
$\mathfrak g$. It is invariant under extension of the base field.
